% ==============================================================================
% T0 Theory / FFGFT: Document 378 (Chapter Content, DE)
% Titel: Begutachtung unter heutigen Bedingungen
% Autor: Johann Pascher
% Datum: 28. September 2026
% Referenzen: Dok. 190, 272, 274, 342, 351, 361, 365, 372, 377;
%   F O R M / CVP v1.0 (Vossen Hg.), CIL (Vossen/Pascher), CCR (Vossen),
%   Universal Constraint / FER (Haskins), SIM2 (Krueger).
% ==============================================================================

\providecommand{\BM}{\textbf{[B]}}
\providecommand{\KM}{\textbf{[K]}}
\providecommand{\SM}{\textbf{[S]}}
\providecommand{\XM}{\textbf{[X]}}
\providecommand{\QM}{\textbf{[Q]}}

\chapter*{Begutachtung unter heutigen Bedingungen}
\addcontentsline{toc}{chapter}{Begutachtung unter heutigen Bedingungen}

{\small\noindent\textbf{Zusammenfassung.}\quad\ignorespaces
Das klassische Peer-Review setzt voraus, dass ein Gutachter eine Arbeit in wenigen
Stunden überblicken kann, dass Autor und Gutachter einer Institution angehören und
dass Texte von Menschen geschrieben werden. Keine dieser Voraussetzungen trägt
heute noch allgemein: Korpora umfassen Hunderte Dokumente, unabhängige Forscher
haben keinen Zugang zu den üblichen Einreichungswegen, und Sprachmodelle erzeugen
plausible Texte, Skripte und Prüfberichte in beliebiger Menge. Dieses Dokument
stellt zusammen, was im Austausch des \emph{Information Physics Institute} (IPI),
in Stefaan Vossens Dot-Theory-Governance und in der eigenen Buchführung des
FFGFT-Korpus bereits an Prüfwerkzeugen entstanden ist, und fragt, was davon für ein
zeitgemäßes Begutachtungs-Regelwerk verwertbar ist. Das Ergebnis ist ein Vorschlag
aus zehn Bausteinen: deklarierte Konstitution vor der Prüfung, ein vom Autor
bereitgestellter Einstieg, prüfbare Objekte statt Texte, vorab festgelegte
Vorhersagen, Rollentrennung, ein fortschreibendes Register, ein geregelter
Umgang mit Sprachmodellen und mit Vereinfachungen, lokalisierte statt globale
Urteile und die ausdrückliche Anerkennung, dass es keinen universellen Prüfer gibt.
Das Dokument enthält keine physikalischen Aussagen; es ist eine methodische
Bestandsaufnahme mit Vorschlagscharakter.
}\medskip

\section*{Zeichenerklärung}
\addcontentsline{toc}{section}{Zeichenerklärung}

\medskip
\begin{tabular}{@{}lp{11cm}@{}}
\toprule
\textbf{Kürzel} & \textbf{Bedeutung} \\
\midrule
{[K]}, {[B]}, {[S]}, {[Q]}, {[X]} & FFGFT-Statusmarker: numerisch bestätigt, algebraisch bewiesen,
  strukturell plausibel mit offener Herleitung, aus externer Quelle übernommen, als falsch erkannt \\
{[SETZUNG]} & ausdrücklich gesetzte Grundannahme, nicht hergeleitet \\
IPI & Information Physics Institute, Mailingliste und Arbeitsgruppen \\
F O R M & \emph{Framework -- Operator -- Residual -- Model}, Grammatik des
  \emph{Constitutional Verification Protocol} (CVP) \\
CIL & \emph{Constitutional Interpretation Layer} (Vossen/Pascher) \\
CCR & \emph{Canonical Constitutional Record} der \emph{Constitutional Physics} (Vossen) \\
FER & Analysemethode von Diana Haskins (Universal Constraint) \\
OAP, FAH & Admissibilitäts-Record und \emph{Framework Admissibility History} der Dot Theory \\
O0, O2 & Originatoren-Deklaration und sekundäres Rendering (Dok.~272) \\
SHA-256 & kryptographischer Hash zur Identität eines Dateiobjekts \\
\bottomrule
\end{tabular}

% ============================================================================
\section{Ausgangslage}
% ============================================================================

\subsection{Was das klassische Verfahren voraussetzt}

Das Begutachtungsverfahren der Fachzeitschriften ist für den Aufsatz gebaut: eine
abgeschlossene Arbeit von zehn bis dreißig Seiten, ein Gutachter aus demselben
Fachgebiet, einige Stunden Lesezeit, ein Urteil in drei Stufen (annehmen, überarbeiten,
ablehnen). Es beruht auf vier stillen Voraussetzungen:
\begin{enumerate}
\item Die Arbeit ist in einer Sitzung überblickbar.
\item Der Gutachter teilt mit dem Autor einen Normalbestand an Methoden, gegen den
  Abweichungen auffallen.
\item Die Zugehörigkeit zu einer Institution filtert vor.
\item Ein Text ist Ergebnis menschlicher Arbeit, sein Umfang also ein grobes Maß für den
  Aufwand dahinter.
\end{enumerate}

\subsection{Was sich geändert hat}

\paragraph{Umfang.} Ein Rahmenwerk wie die FFGFT umfasst inzwischen über
dreihundertsiebzig Dokumente mit eigenen Statusmarkern, einem Korrekturregister
(Dok.~190) und Querbezügen zwischen Geometrie, Galois-Theorie und Teilchenphysik.
Das erschließt sich niemandem in einigen Stunden. Das ist kein Sonderfall: jede
ernsthafte Theorie erreicht diesen Punkt, auch etablierte Programme erschließen sich
nicht an einem Nachmittag.

\paragraph{Zugang.} Unabhängige Forscher haben keine institutionelle Adresse;
Verzeichnisdienste nehmen oft nur auf, was bereits über Verlagsplattformen gelistet
ist. Repositorien (GitHub) und Archive mit DOI (Zenodo) sind öffentlich und
zitierfähig, werden aber vom Begutachtungssystem kaum erfasst.

\paragraph{Sprachmodelle.} Text, Code und scheinbare Prüfberichte lassen sich heute
ohne nennenswerten Aufwand erzeugen. Der Umfang eines Textes misst nicht mehr die
Arbeit dahinter. Im IPI-Austausch sind Fälle aufgetreten, in denen ein Sprachmodell-Prompt
als \enquote{Quellcode} eingereicht wurde, in denen ein Prüfblock seine eigenen
Konstanten ausgab und darunter \enquote{PASS} druckte, und in denen ein zweites
Sprachmodell mit geteiltem Kontext eine frühere Analyse wörtlich wiedergab, statt sie
unabhängig zu prüfen.

\paragraph{Mustersuche.} Aus genügend Zahlen lässt sich jedes Muster herauslesen. Das
Bild der Wunderlampe, an der man reibt, bis etwas erscheint, trifft einen realen
Missbrauch von Rechenleistung. Es trifft aber nicht jede Arbeit, die Zahlen enthält.
Der Unterschied liegt in der Reihenfolge: Bei der Mustersuche stehen die Messwerte am
Anfang, bei einer Herleitung stehen sie am Ende als Prüfstein. In der FFGFT steht
die Geometrie $T^4/\mathbb{Z}_3$ als {[SETZUNG]} fest; daraus folgen
algebraisch $\mathbb{Z}_3$, $\mathrm{GF}(27)$, $D_4$ und Trialität; erst daraus
kommen Zahlen, und erst danach werden sie mit Messwerten verglichen (Dok.~365).
Ein Außenstehender sieht beim ersten Blick nur die Zahl, nicht die Reihenfolge, und
vermutet die Lampe.

\paragraph{Vereinfachung und Vorurteil.} Ein kurzer Einstieg ist möglich und liegt
mehrfach vor. Er wird aber häufig selbst als fehlerhaft gelesen, weil der Leser die
Vereinfachung für die Aussage nimmt. Wer mit der Erwartung kommt, dass eine Arbeit
außerhalb der Institutionen nicht stimmen kann, findet in jeder Vereinfachung den
Beleg dafür. Die Last liegt damit doppelt: beim Autor, der den Einstieg so bauen muss,
dass die Reihenfolge sichtbar ist, und beim Prüfer, der vor einem Fehlerurteil die
vollständige Fassung heranziehen muss.

\subsection{Die Frage}

Die Frage ist nicht, ob Begutachtung nötig ist, sondern wie sie unter diesen
Bedingungen aussehen kann: wer prüft was, an welchem Objekt, in welcher Reihenfolge,
mit welcher Verbindlichkeit, und wer behält die Deutungshoheit worüber.

% ============================================================================
\section{Bestand: was bereits existiert}
% ============================================================================

Die folgenden Werkzeuge sind nicht für ein Begutachtungsverfahren entworfen worden. Sie
sind in der Arbeit selbst entstanden, oft an einem konkreten Streitfall. Gerade deshalb
sind sie erprobt.

\subsection{Die Buchführung des FFGFT-Korpus}

\begin{itemize}
\item \textbf{Statusmarker.} Jede Aussage trägt einen Status: {[K]}, {[B]}, {[S]}, {[Q]},
  {[X]}, {[SETZUNG]}. Der Leser sieht sofort, was bewiesen, was gerechnet, was offen und
  was übernommen ist. Die Marker existierten vor der Korrespondenz zum CCR; spätere
  Terminologie beschreibt, was sie tun (Dok.~351, Methodische Notiz).
\item \textbf{Korrekturregister.} Dok.~190 führt Korrekturen an älteren Dokumenten und das
  Schließen offen geführter Punkte fortlaufend; frühere Fassungen sind als Archiv
  eingefroren. Korrekturen werden nicht still eingearbeitet, sondern verzeichnet.
\item \textbf{Prüfskripte.} Zu jedem neueren Dokument gehört ein Skript mit Assertions,
  das \texttt{PASS} oder \texttt{FAIL} je Aussage ausgibt. Viele dieser Skripte sind
  Nulltests: liefert Zufall dasselbe Ergebnis, wird es nicht gewertet.
\item \textbf{Negativbefunde als Abschnitte.} Dok.~342 erteilt Galois-Produkten ab $k\ge5$
  ausdrücklich keinen Status aus der Produktsuche; Dok.~361 hat einen eigenen Abschnitt
  \enquote{Was nicht übereinstimmt}; Dok.~276 hält ein schwaches Ergebnis als ehrlichen
  Negativbefund fest.
\item \textbf{Leitplanke gegen Exponenten-Numerologie.} Jede Größe ist als $\xi^N$
  schreibbar; eine gemessene Skala in einem Exponenten ist keine Vorhersage (Dok.~190, P35).
\item \textbf{Komparator ist nicht Eingang.} Messwerte (PDG) dienen dem Vergleich, nicht
  der Herleitung. Diese Trennung ist in F O R M als Illustration~001 aufgenommen.
\end{itemize}

\subsection{Die Governance der Dot Theory}

Dot Theory (Vossen) ist nach eigener Auskunft kein konkurrierendes Physikmodell, sondern
ein repräsentationales Meta-Programm. Ihre Governance-Schicht ist in Dok.~272
dargestellt: \emph{Lexicon} (Begriffe), \emph{OAP} (Admissibilität), \emph{Matrix}
(Beziehungen zwischen Rahmenwerken), \emph{FAH} (wie sich Admissibilität ändert).
Tragend sind zwei Unterscheidungen:
\begin{itemize}
\item \emph{Preservation of meaning-making} gegen \emph{allocation of meaning}: festzuhalten,
  wie ein Ergebnis verstanden wurde, ist zulässig; zu entscheiden, wer recht hat, nicht.
\item Originatoren-Deklaration (O0) gegen sekundäres Rendering (O2): die Deutung eines
  Rahmenwerks durch Dritte ist als solche gekennzeichnet; das Korrekturrecht bleibt beim
  Urheber. Records sind dauerhaft und anfechtbar; spätere Einträge können verfeinern,
  bestreiten oder ablösen, das frühere Rendering bleibt Teil der Geschichte.
\end{itemize}

\subsection{Die Constitutional Interpretation Layer}

Die CIL ist von Stefaan Vossen und dem Autor gemeinsam verfasst; die Mitautorschaft ist
auf die CIL-Spezifikation begrenzt und erstreckt sich nicht auf Dot Theory oder den CCR.
Ihr Gegenstand ist nicht die Richtigkeit wissenschaftlicher Aussagen, sondern die
Berechtigung, Schlüsse zu ziehen, zu erben und weiterzugeben. Ausgeschlossen sind
ausdrücklich wissenschaftliche Korrektheit, empirische Gültigkeit, mathematische Wahrheit
und ontologische Festlegung. Aus dem Dialog der beiden Autoren stammen die Prinzipien
\begin{itemize}
\item 3 -- \emph{Asymptotic Humility},
\item 6 -- \emph{Boundary Legitimacy / Evidential Admissibility},
\item 7 -- \emph{Recoverability / Reproducibility},
\item 8 -- \emph{Asserted Constitutional Properties}: eine behauptete Eigenschaft
  (Determinismus, Selbstverwaltung, Herleitungsstatus) bleibt autorale Behauptung, bis sie
  unabhängig rückgewinnbar ist.
\end{itemize}
Weitere Prinzipien der Fassung v1.1 benennen die Vokabelfalle (Übernahme des Vokabulars
belegt keine Übernahme der Methode; Ähnlichkeit von Ergebnis oder Terminologie ist keine
Brücke) und den \emph{constitutional drift} (undeklarierte Ausweitung des Geltungsbereichs).
Vom Autor ergänzt wurde die Bedingung, dass der Geltungsbereich an der Stelle der
Herleitung deklariert sein muss und nicht erst im Streitfall nachträglich behauptet
werden kann; damit wird das Prinzip zu einem prüfbaren Test.

\subsection{F O R M und das Constitutional Verification Protocol}

F O R M (Hg.\ Vossen; Autoren Haskins, Leavitt, Pascher, Quílez Zamora, Vossen; v1.0,
fixiert mit SHA-256 \texttt{bd2a1f2d\ldots e8e90c0e}) gliedert jede Prüfung in
\emph{Framework} (erklärter Unterscheidungsraum), \emph{Operator} (die ausgeführte
Operation), \emph{Residual} (was die Operation übersteht) und \emph{Model} (die
rückgewinnbare Darstellung). Die Besetzung der R-Stelle mit \emph{Residual} statt
\emph{Reality} geht auf einen Vorschlag des Autors zurück: \emph{Reality} hätte eine
ontologische Festlegung eingetragen, die das Protokoll gerade vermeiden will.

Aus derselben Korrespondenz stammt das Argument gegen einen universellen Prüfer. Ein
Beweisprüfer wie Lean funktioniert, weil Axiome, Typsystem und zulässige Operationen
vorher vereinbart wurden; der Prüfer kommt an zweiter Stelle. Das Recht arbeitet ebenso:
vor dem Urteil steht ein vereinbartes, anfechtbares Verfahren. Programmiersprachen liegen
dazwischen und zeigen den Punkt am schärfsten: sie sind stark normiert, und doch liefern
verschiedene Bibliotheken für dasselbe Integral verschiedene Ergebnisse, und
$0{,}1+0{,}2=0{,}3$ ist in IEEE-754-Arithmetik falsch, in exakter rationaler Arithmetik
richtig. Wenn schon dort das Ergebnis vom Rahmen abhängt, kann für heterogene
wissenschaftliche Rahmenwerke kein einzelner Prüfer das Ziel sein, sondern nur ein
Verfahren, das die jeweilige Normierung offenlegt und rückgewinnbar macht.

\subsection{FER}

Die Analysemethode FER von Diana Haskins fragt bei jedem Begriff, was ihm seinen Status
verdient hat, und lässt Begriffe nicht deshalb stärker werden, weil sie gut
zusammenpassen. In einer Anwendung von FER im IPI-Austausch vom 24.--25.~September 2026
entstand daraus eine rekursive Prüffolge: erklärter Rahmen $\to$ Operation $\to$ gewonnene Unterscheidungen $\to$
Darstellung $\to$ nächste Operation $\to$ Rückgewinnungstest $\to$ erhaltene, verlorene
oder neu freigelegte Unterscheidung $\to$ revidierte Darstellung. Drei Fälle werden
getrennt: eine Unterscheidung war verfügbar und ist nicht mehr rückgewinnbar; sie war
nie dargestellt; sie war nicht erkannt, aber die Information war vorhanden und eine
spätere Operation legt sie frei. Die Darstellung ist damit nicht das Ende einer Prüfung,
sondern das nächste Prüfobjekt.

\subsection{Praxis in Audits}

Mehrere Prüfverfahren der letzten Monate haben Regeln hervorgebracht, die über ihren Anlass
hinaus gelten:
\begin{itemize}
\item \textbf{Einfrieren und Hashen.} Der Autor friert sein Objekt ein und nennt den
  SHA-256-Hash; der Prüfer bestätigt den Hash unabhängig am erhaltenen Objekt, bevor er
  inhaltlich prüft. Zwei Dateien mit gleichem Titel, gleichem Datum und gleicher
  Versionsangabe waren im SIM2-Verfahren nur über den Hash zu unterscheiden.
\item \textbf{Der Endstatus wird vergeben, nicht beansprucht.} Eine Einreichung, die sich
  selbst \enquote{ADMISSIBLE} bescheinigt, hat damit keinen Status. Den Status vergibt die
  erneute Prüfung.
\item \textbf{Reihenfolge.} Deklarieren $\to$ einfrieren $\to$ implementieren $\to$ ausführen
  $\to$ unabhängig prüfen $\to$ Status vergeben. Im Block-II-Verfahren (MOTHER-GEA,
  25.--27.~September 2026) kamen Implementierungen mehrfach vor der Spezifikation; genau
  dort traten die Fehler auf.
\item \textbf{Zwei konstitutionelle Invarianten} aus demselben Verfahren: \emph{Löser meldet
  Konvergenz $\neq$ stationärer Punkt} (eine absolute Lösertoleranz ist kein
  Stationaritätstest; nötig sind skalennormierte Residuen, die nach dem Lösen unabhängig
  ausgewertet werden) und \emph{geänderte eingefrorene Eingabe $\neq$ dasselbe Objekt}.
  Ein dritter Befund kam hinzu: dieselben eingefrorenen Eingaben mit zwei verschiedenen
  Potentialfunktionen ergeben zwei verschiedene Objekte; die Spezifikation muss die
  Funktion selbst festlegen, nicht nur ihre Parameter.
\item \textbf{Rollentrennung.} In der CTA-FUSION-Erklärung hat der Autor seine vollständige
  Exposition zum eigenen Korpus erklärt und sich damit von der unabhängigen Prüfung der
  eigenen Treue ausgeschlossen. Im SIM2-Verfahren (Krüger) sind Verwahrer der verdeckten
  Aufgaben, Attestor und Bauer getrennte Rollen mit Interessenerklärung; Koautoren gelten
  nicht als neutral.
\item \textbf{Versiegelte Vorab-Vorhersage.} Die Vorhersage, dass Krügers OP8-Neuro-Operator
  U3 keine Struktur über den Markov-Grenzfall hinaus hinzufügt, wurde vor der Rechnung mit
  Hash hinterlegt und durch das Ergebnis bestätigt.
\item \textbf{Lokalisieren statt Urteilen.} Eine Einschränkung in Rahmenwerk~A begründet oder
  widerlegt nichts über Rahmenwerk~B. Ein Autor darf eine Anfrage nach tieferer
  Offenlegung aus berechtigten Gründen ablehnen; die Folge ist nur, dass die unabhängige
  Rückgewinnbarkeit an der deklarierten Grenze endet. Das lokalisiert den Status, es
  verurteilt nicht.
\end{itemize}

\subsection{Was sich als untauglich erwiesen hat}

\begin{itemize}
\item \textbf{Mehrere Sprachmodell-Gutachter aus demselben Modell.} Verschiedene Rollen
  oder Aufforderungen erhöhen die Abdeckung, nicht die Unabhängigkeit: ein Modell hat einen
  Satz blinder Flecken. Ein zweites Modell mit geteiltem Kontext ist ein Echo, keine
  Replikation.
\item \textbf{Selbstbestätigende Prüfblöcke.} Ein Block, der eigene Konstanten mit sich selbst
  vergleicht, kann nicht fehlschlagen und ist deshalb keine Prüfung.
\item \textbf{Hash über das Ergebnis statt über die Quelle.} Er belegt Reproduzierbarkeit der
  Ausgabe, nicht Identität des Programms und nicht physikalische Richtigkeit.
\item \textbf{Behauptete Eigenschaften.} \enquote{Selbstverwaltend}, \enquote{parameterfrei},
  \enquote{frozen} sind Eigenschaften, die belegt werden müssen; eine einzige angepasste
  Konstante in der Beschreibung genügt, um \enquote{parameterfrei} zu widerlegen.
\end{itemize}

% ============================================================================
\section{Vorschlag: zehn Bausteine}
% ============================================================================

Die folgenden Bausteine sind aus dem Bestand abgeleitet. Sie sind kein fertiges
Regelwerk, sondern ein Vorschlag; ein Regelwerk ist selbst eine Setzung und muss von
denen vereinbart werden, die es anwenden.

\subsection*{B1 -- Deklarierte Konstitution vor jeder Prüfung}
Der Autor legt offen, was gesetzt und was hergeleitet ist, in welcher Reihenfolge die
Schichten aufeinander folgen und welcher Status jeder Aussage zukommt. Für die FFGFT
heißt das: Geometrie $\to$ Algebra $\to$ Zahl $\to$ unabhängiger empirischer Vergleich.
Diese Reihenfolge gehört in jeder Darstellung vor die erste Zahl.

\subsection*{B2 -- Ein Einstieg, den der Autor bereitstellt}
Weil niemand von außen die Zeit für den ganzen Korpus aufbringt, liegt die Last beim
Autor: ein ausgewiesener Einstieg, der in etwa einer Stunde zeigt, worauf alles beruht
und woran es scheitern könnte. Für die FFGFT sind das Dok.~365 (das Fundament in vier
Schichten), Dok.~377 (die Schnittstellen zu anderen Rahmenwerken) und die festen
Vorhersagen, etwa $M_W=80{,}420$~GeV und $\sin^2\theta_{23}\in\{4/9,\,5/9\}$ (Dok.~372).

\subsection*{B3 -- Prüfbare Objekte statt Texte}
Geprüft wird ein eingefrorenes Objekt: Dokument, Skript, Eingabedaten, jeweils mit Hash
über die Quelle. Ein Prüfskript muss scheitern können; es vergleicht gegen einen externen,
vorher festgelegten Datensatz, bricht bei fehlender Identität ab, statt einen Ersatzwert
auszugeben, und enthält Nulltests. Wo eine Funktion konstitutiv ist, wird sie selbst
festgelegt, nicht nur ihre Parameter.

\subsection*{B4 -- Vorab festgelegte Vorhersagen}
Was ein Rahmenwerk wagt, wird vor dem Vergleich mit Daten hinterlegt, mit Datum und Hash.
Ein Rahmenwerk, das nur nachträglich passt, hat nichts gewagt. Der Umfang eines Korpus
ist kein Ersatz für eine einzige scharfe Vorhersage.

\subsection*{B5 -- Rollentrennung und Deutungshoheit}
Der Autor behält die Deutungshoheit über seine Setzungen und deren erklärte Bedeutung
(O0). Er hat keine Hoheit darüber, ob eine erklärte Operation tatsächlich ausgeführt wird
und ob das Ergebnis folgt; das entscheidet die Prüfung. Autor, Prüfer und, wo nötig,
Verwahrer und Attestor sind getrennte Rollen; jede Rolle erklärt ihre Exposition. Die
Deutung eines Rahmenwerks durch Dritte ist als sekundäres Rendering gekennzeichnet.

\subsection*{B6 -- Ein fortschreibendes Register}
Befunde des Prüfers, Antworten des Autors und Korrekturen werden fortlaufend und dauerhaft
verzeichnet, nicht überschrieben. Eine zurückgenommene Aussage bleibt als zurückgenommene
sichtbar. Der Status wird vergeben, nicht beansprucht, und jede Vergabe ist anfechtbar.

\subsection*{B7 -- Sprachmodelle als Werkzeug, nicht als Urteil}
Der Einsatz von Sprachmodellen wird offengelegt. Ihre Ergebnisse werden nicht ungeprüft
übernommen, sondern durch Assertions gegen Rechnung und Quelle abgesichert. Eine
Sprachmodell-Prüfung ersetzt kein unabhängiges Urteil; mehrere Rollen desselben Modells
sind keine Mehrheit. Die Verantwortung für jede Aussage bleibt beim Menschen, der sie
unterzeichnet.

\subsection*{B8 -- Vereinfachung als solche ausweisen}
Vereinfachte Darstellungen sind als solche markiert und nennen die vollständige Fassung.
Ein Einwand gegen die Vereinfachung ist kein Einwand gegen die Herleitung; bevor ein
Prüfer einen Fehler feststellt, zieht er die vollständige Fassung heran.

\subsection*{B9 -- Lokalisieren statt global urteilen}
Das Ergebnis einer Prüfung ist ein Status je Aussage (etwa bestätigt, nicht angenommen,
Kandidat), kein Gesamturteil über ein Rahmenwerk. Eine Grenze der Rückgewinnbarkeit
lokalisiert den Status dessen, was dahinter liegt; sie verurteilt es nicht. Ähnliches
Vokabular begründet keine Brücke zwischen Rahmenwerken.

\subsection*{B10 -- Kein universeller Prüfer}
Jede Prüfung setzt eine vereinbarte Normierung voraus, und diese ist selbst eine Setzung.
Ziel ist nicht ein Prüfer für alle Rahmenwerke, sondern ein Verfahren, das die jeweilige
Normierung offenlegt, ererbte Konventionen (Rechenbibliothek, Einheitensystem) von
rahmenspezifischen Annahmen trennt und beide rückgewinnbar macht.

\subsection{Gestufte Prüftiefe}

Um die Zeitfrage praktisch zu beantworten, lässt sich die Prüfung in drei Stufen
gliedern, die nacheinander und von verschiedenen Personen geleistet werden können:
\begin{center}
\small
\begin{tabular}{@{}p{0.14\linewidth}p{0.40\linewidth}p{0.38\linewidth}@{}}
\toprule
\textbf{Stufe} & \textbf{Frage} & \textbf{Objekt und Aufwand} \\
\midrule
1 Einstieg & Sind Setzungen, Reihenfolge und Statusmarker deklariert? Gibt es vorab festgelegte Vorhersagen? &
  Einstiegsdokumente; etwa eine Stunde \\
2 Objekt & Stimmen Hashes? Laufen die Skripte? Können ihre Prüfungen scheitern? Stimmen Eingaben mit der Spezifikation überein? &
  eingefrorene Skripte und Daten; Stunden \\
3 Tiefe & Folgt eine bestimmte Kette tatsächlich aus den Setzungen? &
  einzelne Herleitung mit ihren Vorgängern; Tage \\
\bottomrule
\end{tabular}
\end{center}
Stufe~1 und~2 verlangen kein Fachwissen im Gegenstand des Rahmenwerks und lassen sich
weitgehend verteilen. Stufe~3 verlangt es und ist nur für einzelne Ketten realistisch.
Ein Urteil einer Stufe gilt nur für diese Stufe.

\subsection{Herkunft der Bausteine}

\begin{center}
\small
\begin{tabular}{@{}p{0.11\linewidth}p{0.42\linewidth}p{0.40\linewidth}@{}}
\toprule
\textbf{Baustein} & \textbf{Herkunft} & \textbf{Verwertbarkeit} \\
\midrule
B1 & FFGFT-Statusmarker; F O R M (Framework) & unmittelbar \\
B2 & FFGFT-Einstiegsdokumente & unmittelbar, Aufwand beim Autor \\
B3 & Audits (Hash, Skripte, Block II) & unmittelbar \\
B4 & versiegelte Vorab-Vorhersage (OP8-Neuro) & unmittelbar \\
B5 & Dot Theory O0/O2; CTA-FUSION; SIM2 & Rollen nur bei mehreren Beteiligten besetzbar \\
B6 & Dok.~190; FAH & unmittelbar \\
B7 & Erfahrung im IPI-Austausch & unmittelbar \\
B8 & Erfahrung mit Einstiegstexten & verlangt Disziplin auf beiden Seiten \\
B9 & CIL (Geltungsbereich, Vokabelfalle); Rahmenwerk~A\,$\neq$\,B & unmittelbar \\
B10 & CVP-Korrespondenz (Normierung vor Prüfer) & als Grundsatz, nicht als Verfahren \\
\bottomrule
\end{tabular}
\end{center}

% ============================================================================
\section{Grenzen und offene Punkte}
% ============================================================================

\begin{itemize}
\item \textbf{Aufwand.} Deklaration, Einfrieren, Hashen und Registerführung kosten Zeit. Das
  Block-II-Verfahren hat gezeigt, dass das Verfahren Fehler zuverlässig sichtbar macht; es hat
  aber auch gezeigt, dass viele Runden nötig werden können, bevor überhaupt ein prüfbares
  Objekt vorliegt. Verfahren kann Inhalt nicht ersetzen.
\item \textbf{Neutrale Prüfer sind knapp.} In einem kleinen Kreis sind fast alle Beteiligten
  Koautoren oder Korrespondenten. Rollentrennung ist dann eher Offenlegung der Nähe als
  echte Unabhängigkeit.
\item \textbf{Anerkennung.} Ein Regelwerk, das innerhalb des IPI vereinbart wird, gilt zunächst
  nur dort. Ob Außenstehende ein nach diesen Regeln geprüftes Objekt anders ansehen als ein
  ungeprüftes, ist offen.
\item \textbf{Vorurteil lässt sich nicht regeln.} Die Bausteine machen Vorurteile sichtbar
  (etwa wenn ein Fehlerurteil ohne Bezug auf die vollständige Fassung ergeht), beseitigen
  sie aber nicht.
\item \textbf{Deutungshoheit gegen Prüfbefund.} Die Trennung in B5 ist im Grundsatz klar, im
  Einzelfall nicht immer: ob eine Aussage zur Setzung gehört oder aus ihr folgen soll, muss
  selbst deklariert sein. Hier greift die Bedingung, dass der Geltungsbereich an der Stelle
  der Herleitung festgelegt sein muss.
\end{itemize}

% ============================================================================
\section{Ergebnis}
% ============================================================================

\begin{enumerate}
\item Das klassische Peer-Review setzt Überblickbarkeit, gemeinsamen Normalbestand,
  institutionellen Vorfilter und menschliche Textproduktion voraus; keine dieser
  Voraussetzungen trägt heute allgemein.
\item Im IPI-Austausch, in der Dot-Theory-Governance, in CIL, F O R M und FER sowie in der
  FFGFT-Buchführung sind Werkzeuge entstanden, die diese Lücken an konkreten Fällen
  geschlossen haben.
\item Daraus lassen sich zehn Bausteine und eine gestufte Prüftiefe ableiten. Die meisten sind
  unmittelbar anwendbar; Rollentrennung und Anerkennung hängen davon ab, wie viele Beteiligte
  es gibt.
\item Das Regelwerk selbst ist eine Setzung. Es gewinnt Geltung nur durch Vereinbarung derer,
  die es anwenden, und es ist wie jede Setzung offen zu deklarieren.
\end{enumerate}

\medskip
\noindent\textbf{Prüfskript:} \texttt{2/python/Dok378\_Skripte/pruef\_378\_verweise.py}
(prüft die zitierten Korpusdokumente, Kennzahlen und Hashes)

\medskip
\noindent\textbf{Register (Dok.~190):} kein Eintrag (neues Dokument, keine Korrektur).

\begin{thebibliography}{99}
\bibitem{dok190} J.~Pascher, \emph{Dok.~190 --- Korrekturregister} (fortlaufend, 2026).
\bibitem{dok272} J.~Pascher, \emph{Dok.~272 --- FFGFT und HLV, eingeordnet nach Dot Theory} (Juni 2026).
\bibitem{dok274} J.~Pascher, \emph{Dok.~274 --- FFGFT im IPI-Feld} (Juni 2026).
\bibitem{dok276} J.~Pascher, \emph{Dok.~276 --- HLV-Homologietest} (2026).
\bibitem{dok342} J.~Pascher, \emph{Dok.~342 --- Faktorisierungsklassen in G(6)} (Sept.~2026).
\bibitem{dok351} J.~Pascher, \emph{Dok.~351 --- PDG-Modellabhängigkeit} (Sept.~2026).
\bibitem{dok361} J.~Pascher, \emph{Dok.~361 --- Kagome-Gitter und FFGFT} (Sept.~2026).
\bibitem{dok365} J.~Pascher, \emph{Dok.~365 --- Das Fundament der FFGFT in vier Schichten} (Sept.~2026).
\bibitem{dok372} J.~Pascher, \emph{Dok.~372 --- Das Matrixelement-Prinzip und seine Reichweite} (Sept.~2026).
\bibitem{dok377} J.~Pascher, \emph{Dok.~377 --- Schnittstellen-Übersicht} (Sept.~2026).
\bibitem{form} S.~Vossen (Hg.), D.~Haskins, J.~D.~Leavitt, J.~Pascher, J.~Quílez Zamora,
  \emph{F O R M --- Constitutional Verification Protocol}, v1.0 (2026),
  DOI 10.5281/zenodo.22875795.
\bibitem{cil} S.~Vossen, J.~Pascher, \emph{Constitutional Interpretation Layer (CIL)}, v1.1 (2026),
  \texttt{dottheory.co.uk}.
\bibitem{ccr} S.~Vossen, \emph{Constitutional Physics --- Canonical Constitutional Record}, Bd.~I--III, v1.0 (2026).
\bibitem{uc} D.~Haskins, \emph{Universal Constraint --- Unified Overview} (2026), DOI 10.5281/zenodo.20257133.
\bibitem{sim2} M.~Krüger, \emph{HLV-LIRO SIM2 v1.3.0} (2026), DOI 10.5281/zenodo.21371481.
\bibitem{ffgft} J.~Pascher, \emph{FFGFT} (2026), DOI 10.5281/zenodo.20117635;
  \texttt{github.com/jpascher/T0-Time-Mass-Duality}.
\end{thebibliography}
